\documentclass[11pt,a4paper]{article}
\usepackage[utf8]{inputenc}
\usepackage[T1]{fontenc}
\title{Simulador geoespacial de radiación solar: documentación de implementación}
\author{TerraPulse AI}
\date{\today}
\begin{document}
\maketitle

\section{Objetivo y estado}
Esta funcionalidad proyecta irradiancia directa incidente y absorbida sobre un terreno a partir de un modelo digital de elevación (DEM) NASADEM. La ROI debe venir del selector compartido del usuario: Chiapas es únicamente un caso de prueba, no el límite geográfico del producto. Una ROI fuera de cobertura NASADEM o de los límites operativos debe rechazarse con una explicación, nunca completar datos faltantes. Los resultados se expresarán en W/m$^2$; no representan temperatura ni una medición de radiación solar. Este documento acompaña la implementación incremental: la adquisición y su caché de granulares están en desarrollo; gradientes, normales y un prototipo de mosaico/reproyección tienen pruebas sintéticas. El mosaico exige cobertura geométrica del ROI y limita píxeles fuente agregados antes de abrir los HGT; los límites y la evidencia por borde siguen siendo insuficientes para uso operativo. Se ha probado localmente una sub-ROI con un ZIP real de una sola tesela, no el ROI completo de Chiapas. El cálculo geométrico de irradiancia directa incidente y absorbida ya tiene pruebas sintéticas; la exportación offline de dos GeoTIFF y un manifiesto es un prototipo, no una entrega científica; su verificación de procedencia, límites y concurrencia y la integración con el visor siguen pendientes.

\section{Contrato de ROI variable y límites actuales}
El pipeline recibe geometría GeoJSON EPSG:4326 de la ROI seleccionada por el usuario; no crea un segundo selector ni toma a Chiapas como constante. El prototipo \texttt{prepare\_dem} valida polígonos, latitudes de cobertura NASADEM y presupuestos y rechaza cobertura ausente. La función pura \texttt{required\_nasadem\_tiles} ya deriva identificadores de teselas requeridas según intersección de área con la geometría y candidatos opcionales para un halo de una tesela; no descarga ni autentica. El adaptador 	exttt{acquire\_roi\_tiles} elige una consulta CMR por cada tesela obligatoria que no tenga una constancia local validada, coteja UR/colección/versión, valida el miembro HGT y guarda un registro local de identidad más hash de su contenido; vuelve a usar únicamente ZIP y registro concordantes. Para un ZIP existente sin registro se necesita una consulta CMR inicial aunque no se descarguen de nuevo los bytes. El registro corrobora metadatos, pero no prueba autenticidad criptográfica del ZIP frente a NASA ni coordinación entre procesos. Las teselas de halo son opcionales y solo se agregan desde la caché validada; no se descargan por defecto. La integración con P4/P5 aún no está implementada. Una región que cruza varias zonas métricas, el antimeridiano o excede el presupuesto no se debe procesar con un bbox engañoso. Rechazarla explícitamente es preferible a inventar cobertura o presentar una irradiancia parcial como resultado completo. El límite de teselas y el presupuesto de píxeles son configurables, pero el valor predeterminado de 30 millones de píxeles fuente admite a lo sumo dos teselas HGT completas de 3601 por 3601 muestras; excederlo detiene \texttt{prepare\_dem} antes de abrir los HGT; \texttt{acquire\_roi\_tiles} comprueba tanto el presupuesto fuente como el mosaico conservador antes de llamar a CMR. Las pruebas sintéticas observaron RED de importación (salida 2), después 10 fallas de comportamiento (salida 1), y GREEN de 11 casos enfocados con 4 advertencias y 139 de la suite con 174 advertencias, ambas salidas 0. Se cubrieron hemisferios, bordes exactos, multigeometría sin teselas de relleno y preparación offline de otra ROI con dos teselas sintéticas. La adquisición pasó RED de importación (salida 2), cinco fallas de comportamiento con el stub (salida 1), y una falla de preflight (salida 1); GREEN de 7 casos enfocados con 8 advertencias y suite de 146 con 182 advertencias, ambas salidas 0. Usó solo ZIP sintéticos y proveedor falso, no una descarga Earthdata ni una prueba Windows. Ni la selección ni el registro local demuestran que un ZIP cacheado sea auténtico frente a un checksum de NASA.

\section{Adquisición implementada hasta ahora}
El módulo \texttt{pipeline/solar/nasadem.py} consulta la colección \texttt{NASADEM\_HGT}, versión \texttt{001}, mediante \texttt{earthaccess.search\_data} con el rectángulo envolvente en longitud/latitud. La descarga se solicita a \texttt{earthaccess.download} y se guarda en un directorio de preparación antes de reutilizar un ZIP en caché. La verificación estructural comprueba la existencia de un miembro HGT, su longitud esperada (3601 por 3601 muestras de dos bytes) y su CRC al leerlo. El prototipo coteja nombres del ZIP y HGT con un identificador CMR simulado, admite \texttt{GranuleUR} dentro de \texttt{umm}, usa un bloqueo entre hilos del mismo proceso y rechaza la caché corrupta sin sobrescribirla. Esos controles no prueban por sí solos cobertura poligonal completa, datum ni coordinación entre procesos. La adquisición no está lista para uso operativo. El usuario informó una prueba autenticada con descarga y validación del granular \texttt{n16w094.zip} y reutilización de caché; la ejecución de red no se reprodujo en esta sesión. Una consulta pública CMR del granular \texttt{G2816786888-LPCLOUD} registró \texttt{umm.GranuleUR=NASADEM\_HGT\_n16w094}, colección \texttt{C2763264762-LPCLOUD} y versión \texttt{001}; no confirma por sí sola la autenticación ni el contenido del ZIP. Las credenciales Earthdata no deben incluirse en el repositorio ni mostrarse al cliente.

La verificación independiente de la adquisición ROI ejecutó una vez 18 pruebas enfocadas con 12 advertencias y 146 de la suite con 182 advertencias (salidas 0), usando únicamente ZIP sintéticos y un proveedor falso. Bloqueó la entrega operativa: el registro local solo refleja el hash del HGT que ya estaba en disco y puede asociarse a un ZIP preexistente sin cotejar un hash remoto NASA; tampoco cierra las carreras entre validación, publicación y retorno de rutas ni demuestra coordinación multiproceso/Windows. El presupuesto de píxeles de salida y memoria de \texttt{prepare\_dem} puede fallar después de la adquisición; un halo no disponible deja derivadas NoData. Las pruebas verdes no prueban resultados científicos generales ni habilitan commit o push del DEM.

\section{Recorte de una tesela completa para una ROI parcial}
Si la región de interés (ROI) ocupa solamente una parte de una tesela NASADEM íntegra, no se reconstruye ni se descarga una ``media tesela'': se utiliza el archivo HGT completo como fuente y se limita la salida a la geometría de la ROI. En el prototipo 
\texttt{pipeline/solar/dem.py}, el procedimiento offline es:
\begin{enumerate}
  \item Validar que la geometría GeoJSON esté en EPSG:4326 y que las teselas suministradas cubran el polígono completo; seleccionar la tesela de la ROI y las vecinas disponibles para el halo. Una ROI que cruce el límite entre teselas requiere las fuentes de ambos lados.
  \item Comprobar presupuestos de teselas, píxeles y memoria estimada antes de abrir los ZIP; verificar identidad, dimensiones y lectura del HGT. No se atribuye a las teselas un valor NoData inventado.
  \item Ensamblar las muestras en una grilla común y comprobar que los postes compartidos de teselas adyacentes coincidan. Recortar al rectángulo de la ROI con un margen de muestras para derivadas; los vecinos ausentes no se sustituyen por ceros.
  \item Proyectar la ROI y la elevación hacia una grilla métrica, remuestrear la elevación y calcular máscaras separadas de píxeles del ROI y de halo válido. La comprobación de cobertura exige soporte completo en las celdas cuyo \emph{centro} está dentro del polígono; una celda de la orilla tocada solamente por la geometría puede permanecer sin dato. La máscara final distingue falta de soporte de una elevación física de cero.
  \item Calcular gradientes y normales únicamente donde la elevación y sus vecinos cardinales son válidos. La interfaz \texttt{compute\_prepared\_terrain} primero combina las máscaras del ROI y del halo como entrada del cálculo métrico y después restringe las cinco capas derivadas y su máscara de validez al ROI. El halo aporta vecinos a las derivadas de la orilla, pero sus normales no se exponen como resultado del ROI. Registrar las teselas que contribuyen a los píxeles finales mediante sus pesos de interpolación; no confundir teselas suministradas con teselas utilizadas.
\end{enumerate}
Las pruebas con HGT sintéticos comprueban el recorte y la preservación de centros de píxel de una ROI interior a una tesela completa; no acreditan todavía una ejecución con NASADEM descargado. Si falta una tesela, el ZIP está incompleto, el ZIP válido no contiene el HGT esperado o contiene uno de identidad distinta, o falta soporte en un centro de píxel del ROI, la preparación emite una advertencia \texttt{RuntimeWarning} y se detiene con \texttt{ValueError}; no entrega un DEM ni irradiancia parcial. Este comportamiento queda a nivel de \texttt{pipeline/solar/dem.py}; no existe todavía una interfaz de usuario o API que muestre la advertencia ni un reintento de descarga verificado.

\section{Orden de las ecuaciones}
El siguiente orden es el \emph{contrato matemático de la funcionalidad}. Los pasos de gradientes y normales tienen pruebas sintéticas y una ejecución acotada con un DEM real; las ecuaciones de irradiancia tienen pruebas sintéticas y una ejecución ilustrativa sobre esa sub-ROI, sin representar radiación medida. Cada nueva fase debe contrastar su implementación con esta secuencia:
\begin{enumerate}
  \item Reproyectar el DEM a coordenadas métricas, interpretar elevación $z$ en metros y propagar la máscara de cobertura/NoData. Nunca calcular gradientes espaciales a partir de grados de latitud/longitud.
  \item Calcular derivadas $z_x=\partial z/\partial x$ y $z_y=\partial z/\partial y$ en m/m, con $x$ hacia el este e $y$ hacia el norte. En una imagen norte-arriba, el índice de fila aumenta hacia el sur; respetar el signo del tamaño de píxel.
  \item Calcular la normal unitaria de cada píxel válido:
  \begin{equation}
    \vec n=\frac{(-z_x,-z_y,1)}{\sqrt{z_x^2+z_y^2+1}}.
  \end{equation}
  Calcular después la pendiente $\arctan\sqrt{z_x^2+z_y^2}$ y la orientación de máxima bajada si se muestran estas capas; ambas derivan del mismo gradiente, no de la irradiancia.
  \item Convertir el azimut $A$ (desde el norte, horario) y la elevación angular $e$ (sobre el horizonte) a radianes y construir el vector hacia la fuente en ejes este, norte y arriba:
  \begin{equation}
    \vec s=(\cos e\sin A,\;\cos e\cos A,\;\sin e).
  \end{equation}
  \item Aplicar la ley del coseno de Lambert a la irradiancia $G_0$ perpendicular al haz, en W/m$^2$:
  \begin{equation}
    G_i=G_0\max(0,\vec n\cdot\vec s).
  \end{equation}
  El valor $G_i$ está en W/m$^2$ de superficie inclinada.
  \item Aplicar el albedo uniforme $\alpha\in[0,1]$ a la irradiancia ya calculada:
  \begin{equation}
    G_{\mathrm{abs}}=(1-\alpha)G_i.
  \end{equation}
  Por tanto, para datos válidos, $0\leq G_{\mathrm{abs}}\leq G_i\leq G_0$. Enmascarar las salidas cuando la normal no sea válida; NoData no es irradiancia cero.
\end{enumerate}
Al cambiar solo $G_0$, $A$, $e$ o $\alpha$ se reutilizarán el DEM y las normales; no se repetirá una adquisición válida. La sombra topográfica, radiación difusa, conducción térmica y temperatura superficial están fuera de alcance. GCOV será una capa radar independiente, nunca un albedo implícito.

\section{Interfaz, entradas y salidas}
El proveedor acepta un rectángulo de búsqueda para identificar candidatos y un identificador de granular para caché; el adaptador de ROI deriva los rectángulos de las teselas obligatorias. \texttt{pipeline/solar/terrain.py} recibe una matriz de elevación, afín GDAL norte-arriba, máscara y unidades explícitas en metros; devuelve gradientes, pendiente en radianes, orientación de máxima bajada en grados, normales y máscara válida. Deriva solo el interior con cuatro vecinos válidos: bordes/NoData quedan enmascarados, las rotaciones y unidades no métricas se rechazan. \texttt{pipeline/solar/dem.py} acepta geometría GeoJSON EPSG:4326 y ZIP HGT; genera una malla proyectada con una máscara de ROI y otra de halo. \texttt{compute\_prepared\_terrain} recibe esa salida, exige CRS proyectado en metros, calcula gradientes, pendiente, orientación y normales mediante \texttt{compute\_terrain} y devuelve \texttt{NaN} fuera de la máscara válida del ROI. En el primer ciclo pasaron nueve pruebas sintéticas, pero una revisión independiente confirmó que aceptaba cobertura incompleta, no limitaba recursos antes de leer teselas ni verificaba costuras. En la continuación siguiente se añadieron casos RED. En las correcciones posteriores se exige cobertura geométrica estricta y se comprueban píxeles del ROI reproyectado, se limita el conjunto de píxeles fuente antes de leer, se cotejan postes compartidos y se densifican bordes para representar la curvatura proyectada. La primera ronda consolidada agregó una estimación de memoria previa a lectura y la segunda preseleccionó teselas. La procedencia se calcula desde pesos de interpolación de las fuentes sobre los píxeles finales del ROI y el halo; la verificación independiente anterior encontró una brecha de cobertura tras el remuestreo y no confirmó la aptitud operativa del DEM. No existe resultado validado con NASADEM real apto para publicación. El adaptador de adquisición por ROI enlaza la selección de teselas con el proveedor, pero su recibo local no demuestra autenticidad de bytes NASADEM ni seguridad multiproceso. La caché derivada L2/L3 experimental queda fuera del flujo operativo. Se implementó un exportador offline prototipo de dos GeoTIFF y manifiesto, bloqueado para publicación por procedencia, concurrencia y recursos; endpoints y visor continúan pendientes de un contrato de integración. El selector de ROI y el backend pertenecen a otros integrantes del equipo y necesitan un contrato conjunto antes de integrarse.

\section{Verificación y límites}
El padre ejecutó \texttt{pytest -q pipeline/tests/test\_solar\_nasadem.py} y observó 11 pruebas aprobadas (código de salida 0). Esas pruebas usaron ZIP sintéticos y mocks de \texttt{earthaccess}; no son una validación de Earthdata, del DEM ni del modelo físico. La corrección cubre en pruebas la forma anidada del resultado CMR e importación sin \texttt{fcntl}; no equivale a una prueba en Windows real. La consulta pública CMR confirmó el patrón de UR de un granular Chiapas, sin descargar bytes en aquella consulta; más tarde el usuario informó la descarga de una tesela, sin que se reprodujera la red aquí; la concurrencia multiproceso, Windows real, reproyección, GeoTIFF y compilación LaTeX siguen pendientes. En el ciclo anterior, \texttt{pytest -q pipeline/tests/test\_solar\_terrain.py} aprobó 7 pruebas; \texttt{pytest -q pipeline/tests/test\_solar\_dem.py} aprobó 9, y la suite aprobó 66. Esos resultados eran anteriores a los nuevos casos de regresión. En esta continuación, el worker ejecutó \texttt{pytest -q pipeline/tests/test\_solar\_dem.py}: salida 1, 5 fallidas y 8 aprobadas (RED). El padre ejecutó \texttt{pytest -q pipeline/tests}: salida 1, 5 fallidas, 65 aprobadas y 28 advertencias. Las fallas observadas son cobertura ROI parcial en dos geometrías, límite de recursos antes de leer, halo y discrepancia en postes compartidos. La verificación independiente, de solo lectura, constató que las pruebas RED no bastan para demostrar presupuesto agregado multitesela, halo por cada borde ni continuidad de una costura válida dentro del ROI; no ejecutó pruebas propias. El plan permite halo válido o NoData explícito donde falten vecinos, sin completar silenciosamente valores. En esa continuación no hubo corrección ni GREEN. En una nueva continuación autorizada, el worker observó RED (6 fallidas, 8 aprobadas) y después GREEN enfocado (19 aprobadas). El padre observó \texttt{pytest -q pipeline/tests/test\_solar\_dem.py} con 19 aprobadas y 33 advertencias, y \texttt{pytest -q pipeline/tests} con 76 aprobadas y 34 advertencias, ambos con salida 0. Aquella verificación independiente bloqueó T3: el umbral geométrico toleraba una franja sin cobertura, no se rechazaban píxeles ROI inválidos tras reproyección, los límites de mosaico y salida se comprobaban después de leer HGT, y las pruebas no demostraban halo/NoData de cada borde ni continuidad de gradiente a ambos lados de una costura. En la nueva continuación autorizada, el worker observó RED (3 fallidas, 22 aprobadas), aplicó una sola corrección y observó GREEN (25 aprobadas). El padre reprodujo \texttt{pytest -q pipeline/tests/test\_solar\_dem.py} (25 aprobadas, 53 advertencias) y \texttt{pytest -q pipeline/tests} (82 aprobadas, 54 advertencias), ambos salida 0. El código ahora exige cobertura geométrica estricta y comprueba píxeles del polígono proyectado, y adelanta presupuestos de mosaico/salida antes de abrir HGT; agregó pruebas sintéticas de halo y costura. La verificación independiente de solo lectura mantuvo T3 bloqueada: la transformación solo de vértices puede omitir bordes curvos del ROI original, el límite por defecto rechaza dos teselas NASADEM completas (25.934.402 píxeles frente a 20.000.000) aunque el ROI sea pequeño, no se acota explícitamente el pico de memoria, se registran teselas no utilizadas como procedencia y la prueba de halo por borde puede aprobar sin observar la derivada del borde. No ejecutó pruebas propias; las cifras GREEN corresponden al padre. En la primera ronda de la continuación consolidada, el worker observó RED (5 fallidas, 24 aprobadas) y GREEN (29 aprobadas); el padre reprodujo \texttt{pytest -q pipeline/tests/test\_solar\_dem.py} con 29 aprobadas y 63 advertencias, y \texttt{pytest -q pipeline/tests} con 86 aprobadas y 64 advertencias, ambos salida 0. El verificador independiente de esa ronda, de solo lectura, mantuvo T3 bloqueada: se leían todos los ZIP antes de seleccionar teselas, la procedencia podía incluir teselas sin aporte, no había límite previo de vértices antes de construir/densificar el ROI, y las pruebas de halo y pendiente en costuras no establecían los resultados específicos esperados. No ejecutó contraejemplos. La segmentación a 0.01 grados tampoco demuestra formalmente un error máximo para todos los ROI, y el presupuesto de memoria es una estimación, no una medición del pico. En la segunda y última ronda de esta continuación, el worker observó RED (3 fallidas, 29 aprobadas) y el padre reprodujo GREEN: 32 pruebas DEM aprobadas con 66 advertencias y 89 de la suite aprobadas con 67 advertencias, ambos exit 0. Se preseleccionan teselas antes de abrir ZIP, la procedencia registra aportes al mosaico y se limita la complejidad de vértices antes de densificar; la prueba de costura contrasta una pendiente analítica. El verificador independiente final, de solo lectura, mantiene T3 bloqueada. La prueba de halo no demuestra la derivada válida con vecino ni NoData cuando falta, para cada borde; la tolerancia de la pendiente analítica en una costura permite incluso el gradiente cero; y la procedencia puede incluir una tesela que solo aportó al recorte fuente, no a píxeles finales del ROI o halo. La densificación no cuenta con una cota de error demostrada y el pico de memoria sigue estimado, no medido; estas últimas son incertidumbres, no contraejemplos ejecutados. Las dos rondas correctivas anteriores concluyeron sin prueba suficiente para declarar completo el DEM. En la nueva unidad enfocada en esas tres garantías, el worker observó RED por procedencia (1 fallida y 31 aprobadas), luego el padre reprodujo GREEN (32 DEM aprobadas, 81 advertencias; 89 de la suite aprobadas, 82 advertencias, salidas 0). Se agregó un mapa de propiedad de muestras reproyectado para identificar teselas presentes en el ROI y el halo finales; las pruebas de bordes distinguen derivadas con y sin tesela vecina, y la prueba de costura exige gradientes analíticos no nulos. La costura ya cumplía la regla antes del cambio, por lo que no se inventa un RED para ella. La verificación independiente, de solo lectura, bloqueó la primera ronda de esta unidad: la elevación bilineal puede mezclar dos teselas mientras que la etiqueta de propiedad por vecino más cercano solo nombra una; las pruebas de halo con y sin vecino no comparan los mismos píxeles del borde. Sí confirmó que la prueba analítica de costura no admite pendiente cero. En la segunda y última ronda, el worker observó RED por soporte de borde (1 fallida, 32 aprobadas); el padre reprodujo GREEN con 33 pruebas DEM aprobadas y 111 advertencias y 90 de la suite aprobadas con 112 advertencias, salidas 0. La preparación exige soporte completo de fuentes para los píxeles válidos y atribuye procedencia a los propietarios con peso positivo en la interpolación bilineal dentro del ROI y halo finales; las pruebas comparan los mismos píxeles de cada borde con y sin tesela vecina. No hubo RED válido previo de la corrección de procedencia bilineal y no se afirma tal evidencia. La verificación independiente final, de solo lectura, bloqueó T3: si un píxel del ROI tiene elevación finita y validez de fuente pero soporte completo insuficiente, el código lo elimina de la máscara válida sin lanzar un error de cobertura. Esto puede aparecer en el borde de una tesela, al recortar parte de ella o junto a una costura; por tanto, las mejoras de halo y procedencia no acreditan cobertura íntegra de la ROI. El presupuesto de memoria es una estimación, no una medición del pico real. El verificador no ejecutó pruebas ni contraejemplos; las 33/90 aprobadas corresponden al padre. Se agotaron las dos rondas de esa unidad y no se hará una tercera automáticamente. Una decisión posterior agregó advertencia y detención ante cobertura o tesela incompleta: el padre observó RED (3 pruebas fallidas y 1 aprobada en los casos enfocados), luego GREEN (5 pruebas enfocadas aprobadas) y la suite \texttt{pytest -q pipeline/tests} con 95 pruebas aprobadas y 124 advertencias. Una prueba adicional confirma que una ROI interior a una tesela completa conserva todos los centros de píxel sin advertencia de incompletitud; la comprobación de cobertura ya no exige soporte de las celdas tocadas solo en el borde y sí falla en centros de píxel incompletos. Una verificación independiente posterior de la cobertura por centros del ROI aprobó 38 pruebas DEM (123 advertencias, salida 0) y constató que el soporte insuficiente ya detiene el cálculo, pero detectó que un ZIP válido sin el HGT esperado o con un HGT de identidad distinta se rechazaba sin \texttt{RuntimeWarning}; por ello T3 no se declaró terminado. En una nueva continuación focalizada, el worker observó RED en los dos casos de identidad (2 fallidas, salida 1); después de una corrección observó GREEN en 40 pruebas DEM (125 advertencias) y la suite de 97 pruebas (126 advertencias), ambas con salida 0 y temporales dentro del proyecto. La verificación independiente de solo lectura confirmó la advertencia antes del rechazo para ambos ZIP y reprodujo 40 pruebas DEM aprobadas con 125 advertencias y la suite de 97 pruebas aprobadas con 126 advertencias, ambas salidas 0 y temporales ignorados dentro del proyecto. No observó RED propio: las dos fallas previas fueron evidencia distinta del worker. Después, el usuario informó autenticación con \texttt{earthaccess}, descarga y validación de \texttt{n16w094.zip} y reutilización de su caché sin una segunda descarga; no se reprodujo ese paso de red en esta sesión. Con el ZIP real ignorado por Git, una ejecución local sin red de \texttt{prepare\_dem} seguida de \texttt{compute\_terrain} sobre una sub-ROI interior de la tesela produjo malla de 361 por 348 píxeles en EPSG:32615, 124022 celdas válidas del ROI y 123625 normales ROI finitas y unitarias (salida 0). Para la nueva composición \texttt{compute\_prepared\_terrain}, el primer intento RED falló durante recolección por ausencia de la interfaz (salida 2); con importación dentro del test se observó RED con 1 falla y salida 1. Después pasaron 8 pruebas de terreno (4 advertencias) y 98 de la suite (130 advertencias), ambas con salida 0. Una ejecución local sin red con la misma tesela real confirmó 123625 normales válidas, finitas y unitarias, sin valores de normales fuera del ROI (salida 0). Esta sub-ROI de una tesela no valida la ROI completa de Chiapas, que cruza también \texttt{n16w093}, ni el pico de memoria para otras regiones, la seguridad de la caché derivada o la coordinación multiproceso/Windows. No se ha declarado DEM operativo ni hubo commit o push. En ese ciclo no se verificaron irradiancia, GeoTIFF, backend o visor. No hay compilador LaTeX disponible: se revisa estructura y comandos, sin afirmar que se compiló un PDF.

\section{Motor matemático de irradiancia directa sin caché derivada}
\texttt{pipeline/solar/radiation.py} implementa \texttt{simulate\_radiation(normals, valid\_mask, irradiance\_w\_m2, azimuth\_deg, elevation\_deg, albedo)} sin importar NASADEM, Django ni la caché L2/L3. Recibe normales unitarias de \texttt{compute\_prepared\_terrain} y aplica, en este orden, el vector solar $\vec s$, la incidencia $G_i=G_0\max(0,\vec n\cdot\vec s)$ y la absorción $G_{\mathrm{abs}}=(1-\alpha)G_i$. Los parámetros finitos cumplen $G_0\geq0$, $A\in[0,360]$ grados, $e\in[0,90]$ grados y $\alpha\in[0,1]$. Rechaza formas incompatibles y normales no unitarias donde la máscara indica datos válidos; fuera de ella devuelve \texttt{NaN}, nunca cero inventado. No recalcula el DEM ni las normales al variar los parámetros. El modelo omite atmósfera, sombras, difusa y temperatura; a elevación cero una ladera puede recibir flujo directo geométrico sin que eso pruebe visibilidad real del Sol.

El test primero falló al importar el módulo ausente (salida 2); con una interfaz mínima sin implementar se observaron 19 fallas de comportamiento (salida 1). Después pasaron 19 pruebas enfocadas y la suite de 125 pruebas con 170 advertencias, ambas salidas 0. Las pruebas cubren el plano horizontal, azimut cardinal y $360^\circ$, albedo cero y uno, irradiancia cero, parámetros inválidos, NoData y conservación de las normales de entrada. Una prueba local sin red encadenó el ZIP real de \texttt{n16w094}, la sub-ROI acotada, normales y el modelo con parámetros \emph{ilustrativos} ($G_0=1000$ W/m$^2$, $A=180^\circ$, $e=45^\circ$, $\alpha=0.2$): 123625 píxeles válidos y $0\leq G_{\mathrm{abs}}\leq G_i\leq G_0$ (salida 0). No es radiación observada ni un resultado publicable de la ROI completa; el exportador GeoTIFF y el manifiesto existen como prototipo offline, pero faltan prueba de procedencia calculada, presupuestos de recursos, validación de cobertura completa e integración visual.

\section{Exportación offline de irradiancia modelada}
\texttt{pipeline/solar/export.py} recibe el resultado incidente y absorbido, el DEM preparado con CRS, afín, máscara y procedencia, los parámetros solares, las huellas SHA-256 del HGT proporcionadas por quien invoca y una versión explícita del código. Valida formas, máscara, valores finitos, cotas físicas y correspondencia $G_{\mathrm{abs}}=(1-\alpha)G_i$ antes de escribir. En un directorio de preparación aislado crea \texttt{incident.tif}, \texttt{absorbed.tif} y \texttt{run.json}; reabre ambos GeoTIFF y verifica sus píxeles válidos, máscara, CRS y transformada antes de exponer el directorio final. El manifiesto registra unidades W/m$^2$, parámetros, identificadores de tesela, hashes HGT proporcionados, versión declarada, CRS, transformada, número de píxeles válidos y hashes de los archivos de salida. La función \emph{no} comprueba por sí misma que las huellas de entrada correspondan a los bytes usados al preparar el DEM: el llamador debe vincularlos; no presenta una medición de radiación real ni usa la caché L2/L3.

Primero falló la importación del módulo nuevo (salida 2); con una interfaz vacía se observaron tres pruebas RED de comportamiento (salida 1). Después pasaron tres pruebas enfocadas y 128 de la suite (170 advertencias), salidas 0. Las pruebas públicas reabren dos archivos y comprueban NoData/georreferencia/manifiesto, rechazan datos incoherentes sin publicar y simulan un error al escribir el segundo GeoTIFF sin dejar una corrida visible. Un ejercicio local sin red con la sub-ROI de \texttt{n16w094.zip} produjo dos GeoTIFF reabiertos con la misma máscara y 123625 píxeles válidos; la versión \texttt{uncommitted-local-solar-v1} y los controles solares fueron solo ilustrativos, no una versión comprometida ni irradiancia observada. La verificación independiente reprodujo tres pruebas enfocadas y 128 de la suite, con salida 0, pero no aprobó una publicación de producción: las huellas HGT y la versión son declaradas por el llamador, no se coteja la incidencia con las normales para los ángulos declarados, no se ha probado la carrera entre bloqueo y renombrado ni la recuperación de un bloqueo abandonado y la serialización completa carece de presupuesto explícito de memoria y disco. El subconjunto de máscara válida debe responder a las normales realmente calculadas, no interpretarse como cobertura ROI completa por sí solo. Falta contrato de directorio con el backend; no se deben distribuir estos archivos como resultados científicos finales.

\section{Caché derivada experimental y límites de entrega}
\texttt{pipeline/solar/preparation.py} ofrece \texttt{prepare\_cached\_terrain}: calcula una clave a partir de la geometría normalizada del ROI, identificadores y huellas de los ZIP fuente, versión espacial, versión de terreno y límites; reutiliza por separado un DEM proyectado (L2) y los campos de terreno (L3). La lectura verifica el artefacto y rechaza la corrupción sin sobrescribirla; la publicación de un archivo usa un enlace duro exclusivo. Esto no constituye una garantía de seguridad entre procesos ni una entrega lista para uso operativo. En pruebas offline sintéticas se observaron ocho casos aprobados con 40 advertencias y la suite de 106 pruebas aprobadas con 170 advertencias, ambas con salida 0; la revisión independiente reprodujo esos resultados. El único RED del worker ocurrió al importar la interfaz inexistente (salida 2); no se registró RED de comportamiento por cada regla. En una prueba local sin red con \texttt{n16w094.zip} y una sub-ROI pequeña, la primera preparación no encontró L2/L3, la segunda reutilizó ambos niveles y conservó 123625 normales válidas (salida 0). Una medición independiente de esa sub-ROI, sin caché derivada, dio 0.221 s, pico RSS del proceso de 153196 KiB y línea de base de 71832 KiB: no es un límite general ni una medición de la serialización de la caché.

La verificación independiente bloqueó la aptitud de L2/L3: entre calcular la huella del ZIP y procesarlo, la fuente podría cambiar y asociar bytes incorrectos con la clave; una ruta de caché con enlace simbólico en un componente intermedio puede aceptarse; y el presupuesto de memoria no incluye las copias al serializar, leer y descomprimir artefactos. El enlace duro tampoco se probó en Windows. Los casos de concurrencia, corruptelas concurrentes y cobertura multitesela permanecen pendientes. No utilizar esta caché para publicar DEM o irradiancia ni hacer commit o push como parte DEM completa hasta corregir y verificar esas garantías. No se completó ninguna descarga adicional ni se modificó el ZIP original.

\section{Preparación temporal y atmosférica de fase A (parcial)}
\texttt{pipeline/solar/time.py} ofrece \texttt{resolve\_solar\_inputs(datetime\_local, timezone\_name, *, altitude\_m, temperature\_c=None, pressure\_hpa=None)}. Recibe fecha-hora local ISO con segundos y sin offset, zona IANA y altitud representativa proporcionada por quien invoca. Resuelve únicamente instantes locales inequívocos: rechaza horas inexistentes y ambiguas por horario de verano, formatos inválidos y zonas desconocidas. Devuelve fecha UTC con offset, offset local, temperatura y presión efectivas junto con su procedencia. No selecciona por sí mismo el punto de la ROI ni comprueba que la altitud proporcionada provenga de NASADEM; por ello no certifica trazabilidad de origen del DEM.

La temperatura se expresa en $^\circ$C, puede proporcionarse y toma 12~$^\circ$C como supuesto predeterminado. La presión se expresa en hPa: puede proporcionarse o estimarse como atmósfera estándar seca a partir de la altitud recibida, con $P_0=101325$~Pa, $T_0=288.15$~K, $L=0.0065$~K/m, $g=9.80665$~m/s$^2$ y $R=287.053$~J/(kg\,K), mediante $P(h)=P_0(1-Lh/T_0)^{g/(RL)}$. Se registran por separado valores proporcionados y predeterminados o calculados; no se presentan como observaciones meteorológicas. Se rechazan valores no numéricos o no finitos, temperatura en o por debajo del cero absoluto, presión no positiva y altitudes fuera del dominio de la expresión cuando se estima la presión. Los controles frontend y la procedencia NASADEM verificada requieren todavía contratos de integración.

La función pura \texttt{resolve\_solar\_inputs} valida fecha y zona, resuelve UTC/offset y calcula presión predeterminada solo cuando falta; \emph{por sí sola} no calcula SPA ni ángulos. La integración posterior en \texttt{pipeline/solar/position.py}, documentada abajo, fija pvlib y selecciona un punto interior a partir de un DEM preparado. No calcula DNI, flujo sobre terreno ni mapas Inferno. La fase B implementa después Ineichen/Linke con climatología; siguen pendientes los límites de la aproximación espacial de un punto, albedo y balance de materiales. Ninguna irradiancia se declara observada.

\subsection{Verificación reproducible}
Desde la raíz del proyecto: \texttt{pytest -q pipeline/tests/test\_solar\_time.py} aprobó 34 pruebas (salida 0); \texttt{pytest -q pipeline/tests} aprobó 205 pruebas con 182 advertencias (salida 0). Los casos añadidos después de una primera corrida detectaron cuatro fallas antes de su corrección: procedencia excesiva y aceptación de números en forma de texto. El RED inicial del worker fue de importación (salida 2), sin RED de comportamiento por cada regla; no se atribuye a esas pruebas validación científica. No se encontró \texttt{pdflatex} en el entorno activo; la compilación del documento queda pendiente.

\section{Fase A: posición solar sobre NASADEM}
La función \texttt{compute\_solar\_position} en \texttt{pipeline/solar/position.py} recibe una ROI GeoJSON EPSG:4326, fecha y hora local con zona IANA, un proveedor con la interfaz existente \texttt{acquire\_roi}, temperatura opcional en $^\circ$C y presión opcional en hPa. No acepta ángulos solares suministrados por la interfaz. Solicita al proveedor las teselas requeridas y convierte sus granulares a \texttt{\{tile\_id: path\}} para \texttt{prepare\_dem}; no implementa una segunda descarga ni usa la caché derivada L2/L3. No se realizó una descarga de red para probar esta fase.

El DEM HGT se prepara sobre grilla métrica y conserva máscara y teselas contribuyentes. Dentro de la ROI proyectada se calcula \texttt{representative\_point} y se elige el centro de píxel de elevación finita y válida que esté \emph{estrictamente dentro} del polígono y minimice su distancia euclidiana métrica a ese punto; los empates se resuelven por fila y columna. Si no existe tal centro, se rechaza el cálculo. Este criterio discreto es reproducible, pero no garantiza representatividad de toda la ROI, de su relieve o de su atmósfera. Se anotan fila, columna, longitud, latitud y elevación del píxel seleccionado. La guía NASADEM HGT v1.3 especifica altura ortométrica en metros relativa al geoide EGM96 y rango válido de enteros $[-32767,32767]$; no se interpreta como altura elipsoidal. El rango físico y la máscara deben controlarse \emph{antes} de usar elevación remuestreada: una muestra fuera de rango no se puede convertir en altitud válida por interpolación. La fuente de la tesela queda identificada, pero un identificador y un hash local no prueban autenticidad frente a una suma oficial de NASA.

Luego se resuelve la hora local a UTC de modo inequívoco. Si no se proporciona presión, se aplica la ecuación de atmósfera estándar seca documentada en la sección previa usando la altura seleccionada; en caso contrario se conserva la presión de entrada. \texttt{pvlib.solarposition.spa\_python} recibe presión $P_{\rm Pa}=100 P_{\rm hPa}$, temperatura efectiva $T_{^\circ\!C}$, altura HGT en metros, latitud/longitud y un \texttt{DatetimeIndex} con zona UTC. La opción \texttt{delta\_t=None} calcula la diferencia temporal a partir de año/mes en la versión fijada; no representa una medición astronómica local. SPA devuelve azimut norte-horario, elevación geométrica $e_g$, elevación aparente $e_a$, cenit geométrico $z_g$ y cenit aparente $z_a$. Se exige $e_a+z_a=90^\circ$ dentro de tolerancia numérica y se identifica noche cuando $e_a\leq 0^\circ$. Se registran UTC, offset, coordenadas, datum, presión/temperatura efectivas y su procedencia, método y versión de pvlib. No hay DNI, radiación modelada, mapas Inferno ni API en esta fase.

Se fija \texttt{pvlib==0.16.1} en \texttt{pipeline/solar/requirements.txt}. Para reproducir pruebas en un entorno que no tenga esta dependencia instalada, usar una instalación aislada del proyecto (sin archivos de datos científicos en Git): \texttt{python -m pip install --target data/solar-pvlib-deps --no-deps -r pipeline/solar/requirements.txt} y luego \texttt{PYTHONPATH=data/solar-pvlib-deps pytest -q pipeline/tests/test\_solar\_position.py}. La opción \texttt{--no-deps} presupone que el entorno ya contiene las dependencias requeridas por pvlib; no es una receta de instalación general. La comparación independiente de software usa el ejemplo A.5 de NREL (17 de octubre de 2003, 12:30:30 hora estándar UTC$-7$, 39.742476$^\circ$N, 105.1786$^\circ$O, 1830.14 m, 820 hPa, 11~$^\circ$C), que publica azimut solar 194.34024$^\circ$ y cenit aparente 50.11162$^\circ$; la coincidencia no valida la atmósfera de otra ROI. Se observó RED de importación (salida 2) antes de crear \texttt{position.py}, no un RED de comportamiento por cada regla. Una prueba posterior mezcló un poste HGT fuera del rango válido ($-32768$) con alturas de 100 m y reveló una altitud interpolada espuria de $-26739.09765625$ m: RED de 3 casos; ahora la lectura HGT advierte y se detiene antes de interpolar muestras fuera de rango. Después de la corrección, \texttt{PYTHONPATH=data/solar-pvlib-deps pytest -q pipeline/tests/test\_solar\_dem.py pipeline/tests/test\_solar\_position.py} aprobó 49 pruebas (175 advertencias) y \texttt{PYTHONPATH=data/solar-pvlib-deps pytest -q pipeline/tests} aprobó 214 (232 advertencias), salidas 0; un verificador independiente reprodujo ambas cifras. Detectó además un comentario obsoleto que confundía el rango válido con un supuesto código NoData y describía el estado anterior del código; se corrigió el comentario sin alterar la lógica. Se leyó sin red una sub-ROI interior de un ZIP NASADEM existente de una tesela y se comprobaron identificador de tesela, píxel seleccionado, datum declarado y $e_a+z_a=90^\circ$ (salida 0); no se verificó autenticidad remota del ZIP ni toda la ROI de Chiapas. La selección de un punto por región no tiene todavía una cota de sensibilidad espacial aprobada ni un benchmark al máximo de píxeles: no usar el resultado para declarar uniformidad, cobertura operacional general ni precisión observacional.

\section{Fase B: DNI directa normal de cielo despejado}
\texttt{pipeline/solar/clearsky.py} ofrece \texttt{compute\_clear\_sky\_dni} y recibe la misma ROI, fecha/hora local, zona y proveedor de fase A, más el porcentaje experimental de intensidad en $[0,100]$ y temperatura/presión opcionales. Valida el porcentaje antes de adquirir datos, llama \emph{una sola vez} a \texttt{compute\_solar\_position} y reutiliza su instante UTC, cénit aparente $z_a$, altitud EGM96, presión efectiva en Pa y temperatura. No repite SPA ni estima presión a partir de la altitud cuando se proporcionó una presión; así el valor del usuario afecta tanto la refracción de fase A como la masa óptica de fase B.

La masa óptica relativa $m_r$ se obtiene con Kasten--Young (1989) sobre $z_a$ y la absoluta adimensional es $m_a=m_r P/101325\,\mathrm{Pa}$, usando la presión efectiva de fase A \cite{kastenyoung-solar}. Se calcula irradiancia extraterrestre $I_0(t)$ en W/m$^2$ con el método Spencer de pvlib. La turbidez Linke $T_L$ es adimensional y procede de \texttt{LinkeTurbidities.h5} empaquetado en pvlib 0.16.1: son \emph{promedios mensuales históricos}, interpolados a día según el día del año \emph{UTC}; no hay turbidez observada para esa fecha ni modelo de nubes \cite{pvlib-linke-solar}. Se registra nombre de LUT, versión del paquete, SHA-256 de la instalación leída, interpolación, $T_L$, $m_r$, $m_a$, $I_0$ y método. La estimación $\mathrm{DNI}_{\rm cs}$ usa \texttt{pvlib.clearsky.ineichen} con $z_a$, $m_a$, $T_L$, altitud, $I_0$ y \texttt{perez\_enhancement=False}; las expresiones GHI, $B_1$, $B_2$ y DNI de la variante fijada están en la sección de cielo despejado de \texttt{docs/modelo-matematico-solar.tex} \cite{ineichen-solar,pvlib-ineichen-solar}. La fase B usa \emph{solo DNI}: no suma GHI ni DHI a la irradiancia de terreno.

Para intensidad experimental $p$ en porcentaje, $\mathrm{DNI}_{p}=\mathrm{DNI}_{\rm cs}\,p/100$ en W/m$^2$; $p$ no es nubosidad, ajuste instrumental ni una probabilidad. Se exige $0\leq\mathrm{DNI}_{p}\leq\mathrm{DNI}_{\rm cs}$ y valores finitos no negativos. Si la elevación aparente $e_a\leq0$, ambas DNI son $0$ W/m$^2$; $T_L$, masas ópticas e $I_0$ quedan \texttt{null} porque no se evalúan, sin inventar datos ni confundir cero físico con falta de una tabla. En día, ausencia de LUT o salidas no finitas/negativas producen error, nunca un reemplazo ficticio. La salida conserva los metadatos de fase A y distingue la estimación climatológica de cualquier observación meteorológica.

\subsection{Comprobación y ejemplo numérico acotado}
Pruebas con pvlib 0.16.1 aislado: \texttt{PYTHONPATH=data/solar-pvlib-deps pytest -q pipeline/tests/test\_solar\_clearsky.py} aprobó 18 casos y \texttt{PYTHONPATH=data/solar-pvlib-deps pytest -q pipeline/tests} aprobó 232 casos (salida 0 del worker). La primera ejecución RED falló en recolección por ausencia del módulo (salida 2), no hubo RED de comportamiento independiente por cada regla. La comparación contra composición directa de las llamadas públicas de pvlib usa entradas sintéticas, no mediciones. Casos de noche, horizonte, porcentajes 0/50/100 y errores de turbidez, DNI o masa no válidas están cubiertos. El padre añadió después un caso de porcentaje entero inmensamente mayor que 100: falló en RED porque se propagaba un \texttt{OverflowError} en lugar de \texttt{ValueError}, y la corrección de la validación dio GREEN en el filtro (8 casos). Se añadió también una prueba pública que encadena A$\to$B con HGT sintético y verifica altitud EGM96, presión, SHA de la LUT y escala de DNI. Cifras finales del padre: \texttt{pytest -q pipeline/tests/test\_solar\_clearsky.py} con 20 aprobadas (35 advertencias) y \texttt{pytest -q pipeline/tests} con 234 aprobadas (267 advertencias), ambas salida 0.

Una \emph{ejecución ilustrativa offline}, sin red ni prueba de autenticidad remota, leyó un ZIP NASADEM local de una sub-ROI interior de una tesela: ROI poligonal acotada $[-93.6,-93.5]\times[16.7,16.8]$ grados, 15 de noviembre de 2026 a las 12:00:00 \texttt{America/Mexico\_City}, $p=50\%$, temperatura predeterminada 12~$^\circ$C y presión estándar estimada 923.0990 hPa desde 779.0168 m EGM96. Se obtuvo $e_a=54.63976^\circ$, $T_L=3.9$, $m_a=1.116264$, $I_0=1397.0985$ W/m$^2$, $\mathrm{DNI}_{\rm cs}=880.7865$ W/m$^2$ y $\mathrm{DNI}_{p}=440.3932$ W/m$^2$ (UTC 18:00:00); el hash de la LUT fue \texttt{34c7e4c0bf865416ecbbfc32b0c3e9b56039a6aadff6997cc3c9ccae8c63e931}. Son \emph{salidas de software condicionadas a entradas y supuestos}, no DNI medida en Chiapas ni validación radiométrica. El punto único, la precisión espacial y la ausencia de nubes reales, sombras y difusa limitan cualquier interpretación; no se implementaron API, Inferno ni balance de materiales.

\section{Fase C: capa de irradiancia directa incidente}
\texttt{pipeline/solar/surface.py} ofrece \texttt{compute\_incident\_irradiance} con la misma entrada de fase B (ROI, fecha/hora local con zona, proveedor, porcentaje, temperatura y presión opcionales). Orquesta una única adquisición de teselas, prepara el DEM métrico, calcula la posición (fase A, sobre el DEM ya preparado, sin repetir adquisición ni SPA), estima la DNI ajustada (fase B, con la posición inyectada) y deriva el terreno con \texttt{compute\_prepared\_terrain}. La proyección usa el motor existente \texttt{simulate\_radiation}: para cada píxel válido con normal unitaria $\vec n=(n_E,n_N,n_U)$ y vector solar unitario $\vec s=(\cos e_a\sin A,\;\cos e_a\cos A,\;\sin e_a)$ en ejes este--norte--arriba, la irradiancia directa incidente por área real de terreno es
\begin{equation}
 G_i(x,y)=\mathrm{DNI}_{p}\,\max\!igl(0,\;\vec n(x,y)\cdot\vec s\bigr)\quad[\mathrm{W/m^2}],
\end{equation}
con $A$ el azimut norte-horario y $e_a$ la elevación aparente de fase A. Es proyección geométrica pura: no supone BRDF Lambertiana, ni añade difusa, sombras, reflejos del terreno o atmósfera píxel a píxel.

Antes de proyectar se verifican consistencias y se falla cerrado ante cualquier anomalía: CRS proyectado en metros; normales finitas, unitarias ($10^{-6}$) y con componente vertical hacia arriba; $A\in[0^\circ,360^\circ]$, $e_a\in[-90^\circ,90^\circ]$; $\|\vec s\|=1$. La máscara válida del ROI gobierna la salida: los píxeles inválidos quedan en NaN (NoData), nunca en cero, y se comprueba $0\leq G_i\leq\mathrm{DNI}_{p}$ en los válidos. De noche ($e_a\leq0^\circ$) o con producto punto negativo, $G_i=0$ W/m$^2$ \emph{físico} sobre píxeles válidos; como \texttt{simulate\_radiation} rechaza elevaciones negativas, los ceros nocturnos se producen directamente y los metadatos conservan la elevación aparente real (negativa) y el indicador \texttt{night}. La capa absorbida y el balance de materiales (absorción, reflexión, transmisión) quedan explícitamente fuera de esta fase, pendientes de la formulación corregida por el físico; no hay API Django ni visualización Inferno.

Para no duplicar cálculos, \texttt{position.py} expone ahora \texttt{acquire\_provider\_tile\_paths} y \texttt{solar\_position\_from\_prepared}, y \texttt{compute\_solar\_position} delega en ellas sin cambiar su comportamiento; \texttt{compute\_clear\_sky\_dni} acepta \texttt{position=} inyectada y solo llama a fase A cuando no se proporciona. Las pruebas existentes de fases A y B no se modificaron y siguen en verde.

\subsection{Comprobación y ejemplo numérico acotado}
Con pvlib 0.16.1 aislado: \texttt{PYTHONPATH=data/solar-pvlib-deps pytest -q pipeline/tests/test\_solar\_surface.py} aprobó 16 casos (salida 0; advertencias de deprecación de pandas/shapely/rasterio ya conocidas) y \texttt{PYTHONPATH=data/solar-pvlib-deps pytest -q pipeline/tests} aprobó 250 casos (salida 0). El worker observó RED de comportamiento (9 fallidas y luego 6 fallidas, salidas 1) antes del GREEN, corrigiendo el \emph{monkeypatch} de referencias y la cobertura del \texttt{PreparedDEM} sintético. Casos numéricos conocidos: terreno plano horizontal $G_i=\mathrm{DNI}_{p}\sin e_a$ en todos los píxeles válidos; plano inclinado $z=a\,y$ comparado contra la normal analítica $(-z_x,-z_y,1)/\|\cdot\|$ y el vector solar de los metadatos; ladera opuesta al Sol con producto punto negativo da exactamente 0; noche con ceros válidos, NaN fuera de máscara y elevación real conservada; contadores prueban una sola llamada a \texttt{acquire\_roi}, a \texttt{spa\_python} y a la consulta Linke por simulación.

Ejecución ilustrativa offline (sin red, ZIP NASADEM local de una tesela, misma sub-ROI de fases A/B): 15 de noviembre de 2026, 12:00 \texttt{America/Mexico\_City}, $p=100\%$, malla 361$\times$348 en EPSG:32615 con 123{,}625 píxeles válidos; $e_a=54.6398^\circ$, $A=180.4759^\circ$, $\mathrm{DNI}_{p}=880.7865$ W/m$^2$. La capa resultante tuvo media 698.9663 W/m$^2$, máximo 880.7541 W/m$^2$ (ladera encarada al Sol, acotado por $\mathrm{DNI}_{p}$), mínimo 0 W/m$^2$ (laderas opuestas) y NaN exactamente fuera de la máscara válida. Son salidas de software condicionadas a supuestos climatológicos y a un punto solar único por ROI, no una medición ni un mapa validado de Chiapas; falta sombra topográfica, difusa y materiales.

\begin{thebibliography}{9}
\bibitem{nasadem-solar} NASA Land Processes Distributed Active Archive Center, \emph{NASADEM User Guide}, versión 1.3, enero de 2025, sección 3, tabla 1.
\bibitem{nrel-spa-solar} I. Reda y A. Andreas, \emph{Solar Position Algorithm for Solar Radiation Applications (Revised)}, NREL/TP-560-34302, ejemplo A.5, 2008.
\bibitem{pvlib-spa-solar} pvlib Python 0.16.1, \emph{solarposition.spa\_python} y \emph{atmosphere.alt2pres}, documentación de la API, 2026.
\bibitem{kastenyoung-solar} F. Kasten y A. T. Young, ``Revised optical air mass tables and approximation formula'', \emph{Applied Optics}, vol.~28, pp.~4735--4738, 1989, doi:10.1364/AO.28.004735.
\bibitem{ineichen-solar} P. Ineichen y R. Perez, ``A new airmass independent formulation for the Linke turbidity coefficient'', \emph{Solar Energy}, vol.~73, pp.~151--157, 2002, doi:10.1016/S0038-092X(02)00045-2.
\bibitem{pvlib-linke-solar} pvlib Python 0.16.1, \emph{clearsky.lookup\_linke\_turbidity}, documentación oficial; tabla histórica mensual \texttt{LinkeTurbidities.h5} empaquetada.
\bibitem{pvlib-ineichen-solar} pvlib Python 0.16.1, \emph{clearsky.ineichen}, \emph{atmosphere.get\_relative\_airmass}, \emph{atmosphere.get\_absolute\_airmass} y \emph{irradiance.get\_extra\_radiation}, documentación oficial.
\end{thebibliography}
\end{document}
